% =========================================================================
% CAPITOLO 5: BUSINESS CASE E CONCLUSIONI (SPOTIFY)
% =========================================================================
\section{Business case finale e conclusioni della Parte I}

La prima parte del presente elaborato ha sviluppato una pipeline completa di Data Science applicata all'analisi di oltre 35.000 tracce musicali estratte da Spotify, combinando pulizia dei dati, modellazione predittiva supervisionata e raggruppamento non supervisionato. L'integrazione di queste metodologie ha permesso di raggiungere conclusioni metodologiche solide e di formulare scenari applicativi di grande valore operativo per l'industria dello streaming musicale.

\subsection{Sintesi dei risultati metodologici}

\begin{itemize}
    \item \textbf{Qualità del dato ed ETL:} La corretta deduplicazione su Track ID (rimozione di 6.428 record duplicati) si è rivelata determinante per eliminare il rischio di data leakage tra set di training e test. La combinazione di SMOTE ed undersampling ha permesso di compensare efficacemente il naturale sbilanciamento delle classi senza introdurre distorsioni geometriche nelle relazioni tra attributi.
    \item \textbf{Classificazione multiclasse:} Su un problema ad elevata complessità (15 generi distinti e sole 13 feature acustiche, senza testo o metadati dell'artista), i modelli Ensemble hanno dimostrato un'eccellente capacità predittiva. Lo \textbf{Stacking Classifier} ha raggiunto un'accuratezza del \textbf{68.90\%} e un Weighted F1 di \textbf{0.6890}, supportato da una \textbf{Macro-Average ROC-AUC pari a 0.94}. L'analisi della matrice di confusione ha attestato che gli errori avvengono quasi esclusivamente tra generi contigui dal punto di vista timbrico e produttivo (Dark Trap, Underground Rap e Trap Metal).
    \item \textbf{Clustering non supervisionato nello spazio PCA:} La selezione ragionata delle 8 feature acustiche continue e la riduzione di dimensionalità mediante \textbf{PCA a 4 componenti principali (73,96\% di varianza spiegata)} hanno permesso di purificare la metrica euclidea da ridondanze e variabili categoriche. Nello spazio PCA a 4 dimensioni, il modello \textbf{K-Means con $k=5$} ha identificato 5 chiari archetipi musicali (\textit{High-Intensity Bass \& Workout, Extended Club \& Instrumental, Acoustic Chill \& Atmospheric, Groove Melodic Urban \& Pop, Fast-Flow Rap \& High-BPM}), validati da un massimo locale del coefficiente di Silhouette pari a \textbf{0,2829}. Parallelamente, l'algoritmo \textbf{DBSCAN ($\epsilon = 0,60, \texttt{min\_samples} = 5$)} ha evidenziato la natura continua del nucleo musicale principale, isolando con precisione un \textbf{7,70\% di tracce outlier (462 brani)} caratterizzate da durate estreme, parlato anomalo o combinazioni timbriche atipiche.
\end{itemize}

\subsection{Ricadute applicative e scenari di business}

I risultati conseguiti aprono la strada a tre concrete applicazioni commerciali:
\begin{enumerate}
    \item \textbf{Generazione di playlist basate sui mood (K-Means):} Le etichette di genere tradizionali risultano spesso rigide o fuorvianti: canzoni formalmente catalogate in generi diversi possono condividere lo stesso impatto emotivo e ritmico. Il partizionamento in 5 macro-mood consente alle piattaforme di streaming di generare raccomandazioni dinamiche per attività quotidiane (\textit{Workout, Focus, Chillout, Party}) basate unicamente sulle caratteristiche acustiche intrinseche della musica.
    \item \textbf{Controllo qualità e anomaly detection automatizzata (DBSCAN):} Con l'upload giornaliero di centinaia di migliaia di brani da parte di etichette indipendenti e utenti, DBSCAN costituisce un potente strumento automatico di controllo qualità, capace di intercettare file corrotti, podcast catalogati erroneamente come musica o registrazioni amatoriali difettose prima dell'indicizzazione editoriale.
    \item \textbf{Valutazione matematica di nuove uscite e demo (Classificazione Stacking):} Il modello predittivo addestrato può essere impiegato dalle case discografiche e dai team di A\&R (Artists and Repertoire) per analizzare le tracce demo appena registrate, identificandone con precisione matematica il posizionamento di genere e il pubblico di riferimento prima del lancio promozionale.
\end{enumerate}

\subsection{Sviluppi futuri}

Come sviluppi futuri per l'analisi acustica, il lavoro potrebbe essere esteso integrando tecniche multimodali: in particolare, l'analisi testuale dei testi delle canzoni (\textit{lyrics}) mediante modelli di Natural Language Processing (BERT/Transformers) e l'elaborazione degli spettrogrammi audio completi attraverso reti neurali convoluzionali (CNN). L'unione di segnali acustici, linguistici e comportamentali (storico di ascolto degli utenti) consentirà di affinare ulteriormente la precisione dei sistemi di raccomandazione musicale di nuova generazione.

